% ============================================================
% Section 61: 肖特基缺陷的综合热力学修正
% ============================================================
\newpage
\section{固体理论：肖特基缺陷的综合热力学修正}
\label{sec:schottky-full-thermo}

\begin{modelbox}[题目：考虑形成能温度依赖、体积功与振动熵的肖特基缺陷浓度]
晶体中产生的肖特基缺陷将给晶体性质带来影响。
\begin{enumerate}
    \item 产生肖特基缺陷时，于格点上脱落的原子在晶体内不断迁移，最后定位于晶体表面，形成了新的一层，这使得晶体体积增大；
    \item 肖特基缺陷的形成能 $u$ 是温度的函数，$u = u_0 - \beta T$，$u_0$ 为 $0\,\mathrm{K}$ 时的形成能，$\beta$ 是常数；
    \item 晶体中空位缺陷附近的原子振动频率发生改变。
\end{enumerate}

考虑了这些影响，试证明肖特基缺陷浓度 $n$ 为：
\begin{equation}
n = N\left(\frac{\omega_E}{\omega'}\right)^{\!3m}
\exp\!\left(\frac{\beta}{k_{\!B}}\right)
\exp\!\left(-\frac{pV_0}{k_{\!B}T}\right)
\exp\!\left(-\frac{u_0}{k_{\!B}T}\right),
\end{equation}
式中 $\omega_E$ 为完整晶体原子振动频率，$\omega'$ 为空位附近 $m$ 个原子的振动频率，$p$ 为压强，$V_0$ 为平均每个原子占据的体积。
\end{modelbox}

\begin{tipbox}[考点快速分析]
\begin{itemize}
    \item \textbf{热力学势的选择}：等温等压过程用\textbf{吉布斯自由能} $G=U-TS+pV$，而非亥姆霍兹自由能 $F=U-TS$
    \item \textbf{形成能的温度依赖}：$u=u_0-\beta T$ 反映晶格膨胀、电子贡献等对形成能的修正，$-T\Delta S$ 项中会多出一个常数因子 $e^{\beta/k_{\!B}}$
    \item \textbf{体积功}：$p\Delta V=pnV_0$，外压抑制缺陷形成（高压下缺陷更少）
    \item \textbf{振动熵的一般处理}：不预先取高温极限，先用完整的爱因斯坦熵公式，最后可退化为高温情形
    \item \textbf{组态熵的两种视角}：$N$ 个格点产生 $n$ 个空位（$W_1=\binom{N}{n}$），或 $N+n$ 个位置容纳 $N$ 个原子（$W_2=\binom{N+n}{n}$）；低浓度下等价
\end{itemize}
\end{tipbox}

\begin{derivationbox}[标准解题过程]
\textbf{Step 1: 吉布斯自由能变化的框架}

肖特基缺陷的形成是等温等压过程，系统的热力学平衡由吉布斯自由能 $G=U-TS+pV$ 的极小值决定。设形成 $n$ 个肖特基缺陷（空位），吉布斯自由能变化为
\begin{equation}
\Delta G = \Delta U - T\Delta S + p\Delta V.
\end{equation}
下面分别计算三项。

\textbf{Step 2: 内能变化 $\Delta U$}

产生一个空位所需能量为 $u=u_0-\beta T$。形成 $n$ 个空位，内能增加
\begin{equation}
\Delta U = nu = n(u_0 - \beta T).
\end{equation}

\textbf{Step 3: 体积变化与压力功 $p\Delta V$}

每个空位相当于将一个原子从内部搬到表面，晶体增加一个原子体积 $V_0$。产生 $n$ 个空位后，总体积增加
\begin{equation}
\Delta V = nV_0,
\end{equation}
相应的压力--体积功为 $p\Delta V = pnV_0$。

\textbf{Step 4: 组态熵 $\Delta S_{\mathrm{conf}}$}

从 $N+n$ 个格点位置（原 $N$ 个格点加上表面新增的 $n$ 个位置）中选出 $n$ 个作为空位，排列方式数为
\begin{equation}
W = \binom{N+n}{n} = \frac{(N+n)!}{n!\,N!}.
\end{equation}
（\textit{注}：这与 $W=\binom{N}{n}$ 在低浓度下等价，因为 $\ln W \approx n\ln(N/n)+n$。）

利用斯特林近似 $\ln x! \approx x\ln x - x$：
\begin{equation}
\Delta S_{\mathrm{conf}} = k_{\!B}\ln W = k_{\!B}\bigl[(N+n)\ln(N+n) - n\ln n - N\ln N\bigr].
\end{equation}

\textbf{Step 5: 振动熵 $\Delta S_{\mathrm{vib}}$（一般温度）}

采用爱因斯坦模型。完整晶体有 $3N$ 个振动模式，频率均为 $\omega_E$。由配分函数，其振动熵为
\begin{align}
S_0 &= -k_{\!B}\sum_{i=1}^{3N}\left[\frac{\hbar\omega_E}{2k_{\!B}T} + \ln\!\left(1-e^{-\hbar\omega_E/k_{\!B}T}\right)\right] \notag\\
&= -\frac{3N\hbar\omega_E}{2T} - 3Nk_{\!B}\ln\!\left(1-e^{-\hbar\omega_E/k_{\!B}T}\right).
\end{align}

产生 $n$ 个空位后，每个空位周围有 $m$ 个原子振动频率变为 $\omega'$。共有 $3nm$ 个模式频率为 $\omega'$，其余 $3(N-nm)$ 个模式频率仍为 $\omega_E$。此时振动熵为
\begin{align}
S &= -\frac{3(N-nm)\hbar\omega_E}{2T} - \frac{3nm\hbar\omega'}{2T} \notag\\
&\quad - 3(N-nm)k_{\!B}\ln\!\left(1-e^{-\hbar\omega_E/k_{\!B}T}\right) \notag\\
&\quad - 3nmk_{\!B}\ln\!\left(1-e^{-\hbar\omega'/k_{\!B}T}\right).
\end{align}

振动熵变 $\Delta S_{\mathrm{vib}} = S - S_0$，化简得
\begin{align}
\Delta S_{\mathrm{vib}} &= -\frac{3nm\hbar(\omega'-\omega_E)}{2T} \notag\\
&\quad + 3nmk_{\!B}\ln\!\left[\frac{1-e^{-\hbar\omega_E/k_{\!B}T}}{1-e^{-\hbar\omega'/k_{\!B}T}}\right].
\label{eq:schottky-vib-entropy-general}
\end{align}

\textit{高温极限}（$\hbar\omega\ll k_{\!B}T$）：$1-e^{-x}\approx x$，对数项退化为
\begin{equation}
\ln\!\left[\frac{1-e^{-\hbar\omega_E/k_{\!B}T}}{1-e^{-\hbar\omega'/k_{\!B}T}}\right]
\approx \ln\!\left(\frac{\omega_E}{\omega'}\right),
\end{equation}
零点能差项 $-3nm\hbar(\omega'-\omega_E)/(2T)$ 相对较小（$\sim\hbar\omega/T$），可忽略。于是
\begin{equation}
\Delta S_{\mathrm{vib}} \approx 3nmk_{\!B}\ln\!\left(\frac{\omega_E}{\omega'}\right).
\label{eq:schottky-vib-entropy-highT}
\end{equation}

\textbf{Step 6: 平衡条件与缺陷浓度}

总熵变 $\Delta S = \Delta S_{\mathrm{conf}} + \Delta S_{\mathrm{vib}}$。将 $\Delta U$、$\Delta S$、$p\Delta V$ 代入 $\Delta G$，并令 $\partial\Delta G/\partial n = 0$。在 $n\ll N$ 近似下，组态熵求导给出
\begin{equation}
\frac{\partial\Delta S_{\mathrm{conf}}}{\partial n}
\approx k_{\!B}\ln\!\left(\frac{N}{n}\right).
\end{equation}

平衡条件为
\begin{equation}
(u_0 - \beta T) + pV_0 - T\left[k_{\!B}\ln\!\left(\frac{N}{n}\right) + 3mk_{\!B}\ln\!\left(\frac{\omega_E}{\omega'}\right)\right] = 0.
\end{equation}
整理得
\begin{equation}
k_{\!B}T\ln\!\left(\frac{N}{n}\right)
= u_0 - \beta T + pV_0 + 3mk_{\!B}T\ln\!\left(\frac{\omega'}{\omega_E}\right).
\end{equation}
两边除以 $k_{\!B}T$ 并取指数：
\begin{align}
\frac{n}{N}
&= \left(\frac{\omega_E}{\omega'}\right)^{\!3m}
\exp\!\left(\frac{\beta}{k_{\!B}}\right)
\exp\!\left(-\frac{pV_0}{k_{\!B}T}\right)
\exp\!\left(-\frac{u_0}{k_{\!B}T}\right).
\end{align}
即
\begin{equation}
\boxed{
n = N\left(\frac{\omega_E}{\omega'}\right)^{\!3m}
\exp\!\left(\frac{\beta}{k_{\!B}}\right)
\exp\!\left(-\frac{pV_0}{k_{\!B}T}\right)
\exp\!\left(-\frac{u_0}{k_{\!B}T}\right).
}
\end{equation}
\textit{证毕。}
\end{derivationbox}

\begin{formulabox}[答案汇总与各项物理意义]
\begin{equation}
n = N\left(\frac{\omega_E}{\omega'}\right)^{\!3m}
\underbrace{e^{\beta/k_{\!B}}}_{\text{形成能温度修正}}
\underbrace{e^{-pV_0/k_{\!B}T}}_{\text{体积功（外压抑制）}}
\underbrace{e^{-u_0/k_{\!B}T}}_{\text{基本热激活因子}}
\end{equation}
\vspace{4pt}
\begin{center}
\begin{tabular}{lll}
\toprule
\textbf{因子} & \textbf{来源} & \textbf{对 $n$ 的影响} \\
\midrule
$e^{-u_0/k_{\!B}T}$ & $0\,\mathrm{K}$ 形成能 & 基本热激活，$T$ 升高 $\Rightarrow$ $n$ 增大 \\
$e^{\beta/k_{\!B}}$ & $u=u_0-\beta T$ 中 $-\beta T$ 项 & 与 $T$ 无关的常数增强因子（$\beta>0$） \\
$e^{-pV_0/k_{\!B}T}$ & $p\Delta V$ 体积功 & 外压抑制缺陷，$p$ 增大 $\Rightarrow$ $n$ 减小 \\
$(\omega_E/\omega')^{3m}$ & 振动频率软化 & $\omega'<\omega_E$ $\Rightarrow$ 因子 $>1$，显著促进缺陷 \\
\bottomrule
\end{tabular}
\end{center}
\end{formulabox}

\begin{insightbox}[物理图像：四种效应对缺陷浓度的``推拉博弈'']]
缺陷浓度的最终表达式是四种热力学因素竞争的结果：
\begin{itemize}
    \item \textbf{推}（促进缺陷）：位形熵、振动熵（$\omega'\!<\!\omega_E$ 时）、形成能中的 $-\beta T$ 项
    \item \textbf{拉}（抑制缺陷）：$u_0$（晶格束缚能）、外压 $p$（体积膨胀要对抗外界）
\end{itemize}

\textbf{温度依赖的层次}：
\begin{itemize}
    \item $e^{-u_0/k_{\!B}T}$：指数敏感，主导温度依赖
    \item $e^{-pV_0/k_{\!B}T}$：同样是指数，但 $pV_0\ll u_0$（常压下通常可忽略）
    \item $e^{\beta/k_{\!B}}$：常数，不随 $T$ 变化，只改变预因子数量级
    \item $(\omega_E/\omega')^{3m}$：常数（若假设 $\omega_E/\omega'$ 不随 $T$ 变），同样是预因子修正
\end{itemize}

\textbf{高压极限}：当 $pV_0\sim u_0$ 时（如地球深部或高压实验），体积功项不可忽略，缺陷浓度被大幅压制。
\end{insightbox}

\begin{thinkbox}[考场速记：缺陷浓度的``万能公式''结构]
各类点缺陷的平衡浓度可统一写成
\begin{equation}
n = N_{\mathrm{eff}}\;\times\;\underbrace{e^{-u_{\mathrm{eff}}/k_{\!B}T}}_{\text{热激活}}
\;\times\;\underbrace{e^{\Delta S_{\mathrm{vib}}/k_{\!B}}}_{\text{振动熵预因子}}
\;\times\;\underbrace{e^{-p\Delta v/k_{\!B}T}}_{\text{压强项}},
\end{equation}
其中 $N_{\mathrm{eff}}$ 为有效位置数（$N$、$\sqrt{NN'}$ 等），$u_{\mathrm{eff}}$ 为有效形成能（可能含 $T$ 线性项）。

\textbf{从简单到复杂的递进}：
\begin{enumerate}
    \item 最简：$n = N e^{-u/k_{\!B}T}$（单质肖特基，仅组态熵）
    \item 加振动：$n = N(\omega/\omega')^{zd} e^{-u/k_{\!B}T}$（sec60）
    \item 加温度依赖的形成能：多一个 $e^{\beta/k_{\!B}}$
    \item 加外压：多一个 $e^{-pV_0/k_{\!B}T}$
    \item 离子晶体成对缺陷：$N\to N$，指数加 $1/2$
\end{enumerate}
\textit{核心思路}：每引入一个新的热力学自由度（振动、压强、温度依赖的形成能），就在自由能中增加相应项，求导后变成浓度表达式中的一个指数或幂次因子。
\end{thinkbox}
